\subsection{An exact logarithmic diagonal and the late-resolution error}
\label{ssec:logdiag}

The following decomposition identifies the principal operator in the actual
CCM basis, including the correction from zero extension at the endpoints.
It requires neither RH nor positivity of the Weil form. All operators in this
subsection use exact assembly; numerical assembly errors are separate.

\begin{proposition}[Bounded logarithmic remainder]\label{prop:logdiag}
Set $c=\lambda^2>1$, $L=\log c$, and let $A_c$ be the self-adjoint operator
of the full Weil form, expressed in the basis $V_j$. Then $A_c=D_L+K_c$,
where
\[
 D_LV_j=d_jV_j,\qquad
 d_j=\tfrac12\log\bigl(\tfrac14+(2\pi j/L)^2\bigr)-\log(2\pi),
 \qquad \operatorname{Dom}(A_c)=\operatorname{Dom}(D_L),
\]
and $K_c$ is bounded self-adjoint with $\|K_c\|\le\mathcal C(c)$, where
\begin{align}
 \mathcal C(c)={}&\pi-\log(1-c^{-1/2})+\gamma+\tfrac\pi2+\log2
 \nonumber\\
 &+\log c+\sqrt c-c^{-1/2}
   +2\sum_{p^k\le c}\frac{\log p}{p^{k/2}}.\label{eq:logdiagconstant}
\end{align}
The same decomposition restricts to either parity sector.
\end{proposition}
\begin{proof}
Transport $V_j$ to $U_j(x)=L^{-1/2}e^{2\pi ijx/L}$ on $(0,L)$ by
\cite[Prop.~3.2]{CCM2025}. For $a>0$, the multiplier
$m_a(t)=\tfrac12\log(a^2+t^2)$ satisfies
\[
 m_a(t)=\log a+\int_0^\infty e^{-as}\frac{1-\cos(ts)}s\,ds.
\]
Subtract its periodic realization from its whole-line realization applied to
zero-extended functions. The local terms cancel and the difference is the
integral operator
\[
 H_{a,L}(x,y)=\frac12\sum_{k\in\mathbb Z\setminus\{0\}}
       \frac{e^{-a|x-y+kL|}}{|x-y+kL|}.
\]
This identity first holds as a quadratic-form identity on trigonometric
polynomials, by truncating the singular integral and taking its limit.
Each of the two adjacent images is bounded in norm by $\pi/2$: reflect one
endpoint and apply weighted Schur to $1/(x+y)$ with weight $x^{-1/2}$,
using $\int_0^\infty u^{-1/2}(1+u)^{-1}\,du=\pi$.
The images $|k|=m+1$, $m\ge1$, together have norm at most $e^{-amL}/m$.
Hence $\|H_{a,L}\|\le\pi-\log(1-e^{-aL})$. Positive values of this kernel
are not being used to assert positive semidefiniteness.

The archimedean multiplier is
$h(t)=\operatorname{Re}\psi(1/4+it/2)-\log\pi$
\cite[eq.~(3.19)]{CCM2025}. Choose $a=1/2$. With $z=1/4+it/2$ its remainder
after subtracting $m_{1/2}(t)-\log(2\pi)$ is
$\operatorname{Re}(\psi(z)-\log z)$. The digamma integral
\cite[eq.~5.9.13]{DLMF} gives
\[
 \psi(z)-\log z=-\int_0^\infty
   \left(\frac1{1-e^{-s}}-\frac1s\right)e^{-zs}\,ds.
\]
The parenthesis is positive, so the absolute value of the real remainder
is at most $\log(1/4)-\psi(1/4)=\gamma+\pi/2+\log2$.
Plancherel gives the same norm bound after compression to the interval.
The pole kernel $2\cosh((x-y)/2)$ has eigenvalues
$L\pm2\sinh(L/2)$ and norm $L+2\sinh(L/2)$. Each active prime power is
$-(\log p)p^{-k/2}$ times a compressed translation and its adjoint, with
norm at most $2(\log p)p^{-k/2}$. These estimates give
\eqref{eq:logdiagconstant}.

Trigonometric polynomials are a form core for $D_L$ and for the Weil form
\cite[Prop.~3.4]{CCM2025}. A bounded perturbation preserves the form domain;
equality on the core therefore identifies the closed forms and their
associated operators. Reflection about $L/2$ commutes with every component.
\end{proof}

\begin{proposition}[Quantitative tail stability at fixed cutoff]
\label{prop:tailstability}
Suppose $A_c$ has a simple normalized ground state $u$, eigenvalue
$\varepsilon_\infty$, and spectral gap $g>0$. Let $u_N$ be a normalized
ground state of $P_NA_cP_N$, phase aligned with $u$, and put
$q_N=\|(I-P_N)u\|_2^2$. For all sufficiently large $N$,
\begin{equation}
 \|u_N-u\|_2^2\le2\left(1+\frac{\mathcal C(c)^2}{g^2}\right)q_N.
 \label{eq:tailstability}
\end{equation}
For centered logarithmic coordinates and every complex $z$,
\begin{equation}
 |\widehat u_N(z)-\widehat u(z)|
 \le\sqrt{2L\left(1+\mathcal C(c)^2/g^2\right)q_N}\,
       e^{L|\operatorname{Im}z|/2}.\label{eq:tailtransformexplicit}
\end{equation}
\end{proposition}
\begin{proof}
Write $p=P_Nu$, $w=(I-P_N)u$. The projected eigenvalue equation is
$(P_NA_cP_N-\varepsilon_\infty)p=-P_NK_cw$.
By min--max, every retained eigenvalue other than the first is at least
$\varepsilon_\infty+g$. Projecting this equation orthogonally to $u_N$
therefore gives $\|(I-|u_N\rangle\langle u_N|)p\|_2
\le\mathcal C(c)\sqrt{q_N}/g$. Consequently
$1-|\langle u,u_N\rangle|^2\le(1+\mathcal C(c)^2/g^2)q_N$.
Phase alignment and $2(1-t)\le2(1-t^2)$ for $0\le t\le1$ prove
\eqref{eq:tailstability}. Cauchy--Schwarz in the Fourier integral proves
\eqref{eq:tailtransformexplicit}. The ground Ritz value eventually lies
below $\varepsilon_\infty+g$ by form-core convergence, making it simple.
\end{proof}

\begin{proposition}[Late eigenvalue error as logarithmic tail energy]
\label{prop:lateenergy}
Under the preceding spectral hypotheses, let
$E_N=\sum_{|j|>N}d_j|u_j|^2$, $v=P_Nu/\sqrt{1-q_N}$,
$\delta_N=\langle v,A_cv\rangle-\varepsilon_\infty$, and
$B=\mathcal C(c)+|\varepsilon_\infty|$.
If $q_N<1$ and $\delta_N<g$, then
\begin{align}
 \frac{E_N-Bq_N}{1-q_N}
 -\frac{\mathcal C(c)^2q_N}{(1-q_N)(g-\delta_N)}
 &\le\varepsilon_N-\varepsilon_\infty
 \le\frac{E_N+Bq_N}{1-q_N}.\label{eq:lateenergybounds}
\end{align}
If $E_N>0$, $d_{N+1}>0$, and $\delta_N\le g/2$, this implies
\begin{equation}
 \left|\frac{\varepsilon_N-\varepsilon_\infty}{E_N}-1\right|
 \le\frac{q_N}{1-q_N}
 +\frac{B+2\mathcal C(c)^2/g}{(1-q_N)d_{N+1}}.
 \label{eq:lateenergyrelative}
\end{equation}
In particular $\varepsilon_N-\varepsilon_\infty\sim E_N$ at fixed $c$
unless the eigenvector has finite support, in which case convergence terminates.
\end{proposition}
\begin{proof}
The exact projection identity is
$\delta_N=\langle w,(A_c-\varepsilon_\infty)w\rangle/(1-q_N)$;
its difference from $E_N/(1-q_N)$ is bounded by $Bq_N/(1-q_N)$.
Rayleigh--Ritz gives the upper bound. The retained residual of $v$ is
$-(I-|v\rangle\langle v|)P_NK_cw/\sqrt{1-q_N}$, of squared norm at most
$\mathcal C(c)^2q_N/(1-q_N)$. The retained second eigenvalue is at least
$\varepsilon_\infty+g$, so the Temple residual inequality bounds
$\langle v,A_cv\rangle-\varepsilon_N$ by that squared norm divided by
$g-\delta_N$. This proves the lower bound. Finally $E_N\ge d_{N+1}q_N$
gives \eqref{eq:lateenergyrelative}. Since $u\in\operatorname{Dom}(D_L)$,
$E_N,q_N\to0$, while $d_{N+1}\to\infty$.
\end{proof}

These results supply a quantitative fixed-cutoff Galerkin transfer, but their
constants can grow rapidly as the spectral gap closes. They neither establish
a coefficient asymptotic nor control the coupled $c,N$ limit. For example,
if one could prove $|u_j|^2\sim A(c)^2/[j^2(\log|j|)^3]$ for an even state,
then \eqref{eq:lateenergyrelative} would imply
$\varepsilon_N-\varepsilon_\infty\sim2A(c)^2/[N(\log N)^2]$.
The coefficient premise and its nonzero amplitude remain unproved; this is
not a replacement fitted law for the large-cutoff ramp.
